
\documentclass[11pt]{article}
\usepackage{amsmath,amssymb,amsthm,mathtools}
\usepackage[margin=1in]{geometry}

\newtheorem{theorem}{Theorem}
\newtheorem{lemma}{Lemma}
\newtheorem{proposition}{Proposition}
\newtheorem{definition}{Definition}

\title{Rank-Two Affine S-Unit Prime Production for Erdős--Graham 203}
\author{Jared Wilder}
\date{\today}

\begin{document}
\maketitle

\begin{abstract}
This draft gives the intended analytic proof architecture for the rank-two
affine S-unit prime-production statement equivalent to Erdős--Graham 203:
for every \(m\) with \((m,6)=1\), at least one integer \(m2^k3^\ell+1\) is
prime.  The elementary algebraic reductions are explicit; the analytic package
is split into six finite components.
\end{abstract}

\section{Main theorem}

\begin{theorem}[Rank-two affine S-unit prime production]
For every \(m\ge 1\) with \((m,6)=1\), there exist \(k,\ell\ge0\) such that
\[
m2^k3^\ell+1
\]
is prime.
\end{theorem}

\section{Elementary reductions}

\begin{lemma}[Prime-divisor congruence]
If \(p\mid mA+1\), then \(mA\equiv-1\pmod p\).  If \(p\) is prime, then
\((p,m)=1\).
\end{lemma}

\begin{lemma}[Pair-GCD difference]
If \(d\mid mA+1\), \(d\mid mB+1\), and \((d,m)=1\), then \(d\mid A-B\).
\end{lemma}

Taking \(A=2^k3^\ell\) and \(B=2^{k'}3^{\ell'}\), all Type-II common-divisor
terms reduce to divisors of
\[
2^k3^\ell-2^{k'}3^{\ell'}.
\]

\section{Collision identity}

Let \(T_D=\{(k,\ell):k+\ell\le D\}\).  Define
\[
C(D,q)=\#\{(x,y)\in T_D^2:2^{x_1}3^{x_2}\equiv2^{y_1}3^{y_2}\pmod q\}.
\]
For \(h=(a,b)\ne(0,0)\), define
\[
R_{a,b}=\left|2^{a_+}3^{b_+}-2^{a_-}3^{b_-}\right|.
\]
Then
\[
C(D,q)-|T_D|
=
\sum_{(a,b)\ne(0,0)}N_D(a,b)\mathbf 1_{q\mid R_{a,b}},
\]
where
\[
N_D(a,b)=
\begin{cases}
0,&S<0,\\
(S+1)(S+2)/2,&S\ge0,
\end{cases}
\quad
S=\min(D,D-a-b)-\max(0,-a)-\max(0,-b).
\]

\section{Six analytic components}

\subsection{Vaughan--Heath-Brown}
Decompose
\[
\sum_{(k,\ell)\in T_D}\Lambda(m2^k3^\ell+1)
\]
into Type-I, Type-II, and acceptable remainder terms.

\subsection{Type-I distribution}
Use the congruence \(m2^k3^\ell\equiv-1\pmod p\) and distribution of the
rank-two orbit in residue classes.

\subsection{Type-II S-unit difference estimate}
Bound off-diagonal bilinear terms through the divisor sums over \(R_{a,b}\).

\subsection{Short-relation lattice control}
Split primes \(q\) for which a short relation vector \((a,b)\) satisfies
\(q\mid R_{a,b}\).  Use the exact formula for \(N_D(a,b)\).

\subsection{Subgroup concentration}
For \(H_p=|\langle 2,3\rangle\subset\mathbb F_p^\times|\), prove a uniform
estimate of the shape
\[
\sum_{p\le z}\frac1{H_p}\ll \log\log z+O(1).
\]

\subsection{Brun--Hooley parity breaking}
Use the above estimates in a lower-bound sieve to prove
\[
\#\{(k,\ell)\in T_D:m2^k3^\ell+1\text{ prime}\}>0.
\]

\section{Lean correspondence}

The Lean theorem corresponding to the analytic package is:
\[
\texttt{EG203.Analytic.FullAnalyticPackage}.
\]
The final file \texttt{EG203/Analytic/ProveIt.lean} proves EG203 once this
package is supplied.
\end{document}
